\begin{theorem}[Fixed physical route and commutator-class measurement]%
    \label{thm:peps_torus_swept_physical_route_measurement}
    \lean{TNLean.PEPS.IsGIsometric.exists_torusSweptPhysicalRoute_commutatorMeasurement}
    \leanok
    \uses{thm:peps_torus_swept_physical_route_composition,
        thm:peps_torus_swept_route_joint_flux_measurement}
    Assume torus periods at least eight and seven, and regular
    G-isometry of every original site tensor. There are twelve global
    unitary factors, each supported on its actual six-site movement
    region, and a complete projective measurement on the six reunion
    spins, all chosen before both initial flux labels $g,h$.
    Let $U=M_{11}\cdots M_0$ be their chronological product and let
    $\Psi_0(g,h)$ be the literal four-endpoint initial state. Then $U$
    sends it to the actual final route state. If $Q_C$ is the reunion
    projector for conjugacy class $C$ and $\widehat Q_C$ its identity
    extension to the whole physical space, then
    \begin{align}
        \widehat Q_C U\Psi_0(g,h)
        &=\mathbf1_{C=[ghg^{-1}h^{-1}]}U\Psi_0(g,h),\\
        Q_C\mathcal C_R(U\Psi_0(g,h))
        &=\mathbf1_{C=[ghg^{-1}h^{-1}]}\mathcal C_R(U\Psi_0(g,h)).
    \end{align}
    Here $\mathcal C_R$ is the actual coefficient matrix across the
    reunion cut. The complementary physical-range outcome annihilates
    these states. The measured class is the identity precisely when
    $g$ and $h$ commute.
    This is the explicit finite-torus route and subsequent class
    measurement for \cite[\S6.6, local lines 2340--2423]{Schuch2010PEPS}.
    It does not assert an unrestricted braid or partner annihilation.
\end{theorem}

\begin{proof}\leanok
    Choose the twelve elementary physical operations and the reunion
    projectors independently of the flux labels. Their checked
    chronological action identifies the output with the actual final
    contraction. Apply the reunion cut identity to that contraction.
    Evaluation at every assembled region/complement configuration
    gives the whole-state action of the identity-extended projector.
\end{proof}
